\section{多项式}\label{sec:polynomials}

\setlength{\intextsep}{3bp}
\setlength{\columnsep}{1em}
\begin{wrapfigure}[7]{r}{.247\textwidth}
  \centering
  \includegraphics[width=\linewidth]{figures/unit-circle-vector.pdf}
\end{wrapfigure}
多项式及其解是我们学习数学时最早接触的一些对象. 例如, 你或许还记得第一次得知, 在 $\mathbb{R}^{2}$ 中, 二元多项式方程
\[
  x^{2}+y^{2}-1=0
\]
的解构成单位圆. 不过, 如果我们的目标是研究单位圆, 还可以选用许多其他多项式来描述它. 例如, 下列任一方程的解集也都是单位圆:
\[
  2x^{2}+2y^{2}-2=0
  \quad\text{或}\quad
  (x^{2}+y^{2}-1)^{2}=0.
\]
这就自然引出一个问题: 有没有一个描述单位圆的\textbf{最佳}多项式?
代数几何给出的答案, 在某种意义上是最平等的: 所有描述单位圆的多项式都同等重要, 我们应当把它们作为一个\textbf{集合}来研究.
那么, 研究一个多项式\textbf{集合}是什么意思呢?

我们很快就会看到, 描述单位圆的多项式所成的集合远不止是一个集合; 它具有重要的代数结构, 反映了圆的几何性质.
为了能描述这个例子及其他情形中的代数结构, 我们必须先为多项式集合及其代数结构确立准确的记号和术语.
我们从单项式的概念入手, 正式开始讨论多项式.

\begin{Definition}{单项式}\label{def:monomial}
  变量 $x_{1},\ldots,x_{n}$ 的\textbf{单项式}是形如
  \[
    x^{\alpha}=x_{1}^{\alpha_{1}}\cdots x_{n}^{\alpha_{n}},
  \]
  的表达式, 其中 $\alpha=(\alpha_{1},\ldots,\alpha_{n})\in\mathbb{N}^{n}$ 是\textbf{指数向量}.
  两个单项式相等, 当且仅当它们的指数向量相同.
\end{Definition}

\begin{BookFlowNote}{%
  变量只有几个时, 通常用不同的字母表示它们, 例如 \mbox{$x$, $y$ 和 $z$}.
}
变量 $x_{1},\ldots,x_{n}$ 应看作形式符号, 它们的作用只是充当 $\alpha_{1},\ldots,\alpha_{n}$ 的占位符.
一个关于 $x_{1},\ldots,x_{n}$ 的单项式所包含的信息, 与一个指数向量 $(\alpha_{1},\ldots,\alpha_{n})$ 所包含的信息是等价的; 不过, 把每个 $\alpha_{i}$ 写成 $x_{i}$ 的指数, 在下面定义单项式的乘法时会很有用.
指数为 $0$ 的变量通常省略不写. 例如,
\[
  x^{2}y^{3}z^{0}=x^{2}y^{3}
  \quad\text{且}\quad
  x^{0}y^{0}z^{0}=1.
\]
\end{BookFlowNote}
与 $f(x,y)=x^{2}+y^{2}-1$ 一样, 多项式由单项式的线性组合构成.
在最一般的情形下, 这些线性组合的系数属于任意环 $R$, 如下一个定义所示.
